\documentclass[11pt]{article}

\usepackage[margin=1in]{geometry}
\usepackage{times}
\usepackage{amsmath,amssymb}
\usepackage{booktabs}
\usepackage{array}
\usepackage{xcolor}
\usepackage{listings}
\usepackage{enumitem}
\usepackage{hyperref}
\usepackage{caption}
\usepackage{fancyhdr}
\usepackage{titlesec}

\hypersetup{
  colorlinks=true,
  linkcolor=black,
  citecolor=black,
  urlcolor=blue,
  pdftitle={Cryptanalysis of KOD8},
  pdfauthor={ Fatih ULUSOY}
}

\definecolor{codebg}{RGB}{245,245,245}
\definecolor{codekw}{RGB}{0,90,150}
\definecolor{codestr}{RGB}{140,20,20}
\definecolor{codecomm}{RGB}{110,110,110}

\lstset{
  backgroundcolor=\color{codebg},
  basicstyle=\ttfamily\small,
  keywordstyle=\color{codekw}\bfseries,
  commentstyle=\color{codecomm}\itshape,
  stringstyle=\color{codestr},
  showstringspaces=false,
  breaklines=true,
  frame=single,
  framesep=4pt,
  language=Python,
  numbers=left,
  numberstyle=\tiny\color{gray},
  numbersep=6pt,
  xleftmargin=1.2em,
  tabsize=4
}

\pagestyle{fancy}
\fancyhf{}
\fancyhead[L]{\small Cryptanalysis of KOD8}
\fancyhead[R]{\small \thepage}
\renewcommand{\headrulewidth}{0.3pt}

\titleformat{\section}{\normalfont\Large\bfseries}{\thesection}{1em}{}
\titleformat{\subsection}{\normalfont\large\bfseries}{\thesubsection}{1em}{}

\title{\textbf{Structural and Empirical Cryptanalysis of KOD8:\\
A Custom Eight-Stage Symmetric Cipher}}
\author{ Fatih ULUSOY ULUSOY}
\date{August 2026}

\begin{document}

\maketitle

\begin{abstract}
\noindent
KOD8 is a custom symmetric cipher that routes input through one of six
input-type-specific ``CipherLists,'' each composing exactly eight primitive
reversible transformations (encodings, XOR variants, classical
substitutions, and transpositions) driven by a single eight-digit numeric
key. This report presents a structural and empirical cryptanalysis of the
reference implementation. We show that the cipher's advertised eight-layer
design provides far less effective key-dependent entropy than its layer
count suggests: several primitives ignore the key entirely, one collapses
an eight-digit (nominally $10^{8}$) key into at most 325 distinct internal
states, and the block-cipher-inspired CBC-style layer reuses the literal
key string as a fixed, non-random initialization vector across every
message. We further show that the cipher is deterministic (identical
plaintexts always yield identical ciphertexts), unauthenticated (no
integrity tag, with several primitives silently absorbing decoding
errors), and that the reference application persists the selected
CipherList identifier as a cleartext header, eliminating any uncertainty
an attacker would otherwise face in classifying the ciphertext. A
known-plaintext brute-force attack against the full nominal keyspace is
implemented and benchmarked, recovering the key from a single
plaintext/ciphertext pair for the plain-text CipherList in under 2.4 hours
on a single CPU core using an unoptimized reference Python implementation
— and in effectively constant time once the structural entropy collapses
identified in this report are exploited. We conclude that KOD8 is
correctly classified as a composite classical (product) cipher and a
functional obfuscation layer, not as a cryptographically secure primitive,
and we outline concrete remediations.
\end{abstract}

\tableofcontents

% =====================================================================
\section{Introduction}
% =====================================================================

KOD8 (``Kod'' — Turkish for ``code'') is a hand-designed symmetric cipher
implemented in \texttt{kod8\_engine.py}, with a companion Tkinter
front end, \texttt{kod8\_ui.py}, that visualizes the encryption and
decryption process step by step. The engine's own documentation frames it
as an ``anti-LLM and decoding agent based cipher model'' — that is, a
scheme designed primarily to resist casual or automated pattern-based
decoding rather than to meet a formal security definition such as
IND-CPA or IND-CCA2.

This framing matters for how the design should be evaluated. KOD8 is not
positioned as a competitor to AES-GCM or ChaCha20-Poly1305; it is a
personal cryptographic-engineering exercise. Evaluated on those terms it
is a legitimate and instructive artifact: a working example of a
\emph{product cipher} in the classical sense — a chain of substitution and
transposition steps, in the tradition of ciphers that predate the
computational security paradigm.

The purpose of this report is to evaluate that claim rigorously. We ask
three questions:

\begin{enumerate}[leftmargin=1.6em]
  \item How much of the nominal $10^{8}$-key\-space entropy actually
        reaches each of the eight steps in a given CipherList?
  \item Under a realistic threat model for how the reference
        implementation is used (a hardcoded module-level key, and a file
        format that stores the CipherList identifier in cleartext), what
        does an attacker actually have to solve?
  \item Empirically, how long does key recovery take against the
        reference implementation?
\end{enumerate}

All findings below are derived directly from the reference source
(\texttt{kod8\_engine.py}, \texttt{kod8\_ui.py}) and validated by
executing the actual engine code rather than by inspection alone.

% =====================================================================
\section{System Overview}
% =====================================================================

KOD8 defines a single global key,

\begin{lstlisting}
KOD8_KEY = "86247931"
\end{lstlisting}

\sloppy
\noindent
an eight-character decimal string, and a fixed 29-letter Turkish alphabet
used by two of its substitution primitives. Seventeen primitive
operations are defined in total — two encodings
(\texttt{op\_hex}, \texttt{op\_base64}), eight XOR/substitution
variants (\texttt{op\_xor\_key}, \texttt{op\_rolling\_xor},
\texttt{op\_keystream\_xor}, \texttt{op\_vigenere\_tr},
\texttt{op\_atbash\_tr}, \texttt{op\_unicode\_shift}, \texttt{op\_sbox},
\texttt{op\_base36}), and seven transposition/block operations
(\texttt{op\_rail\_fence}, \texttt{op\_block\_rotate},
\texttt{op\_split\_reverse}, \texttt{op\_full\_reverse},
\texttt{op\_columnar}, \texttt{op\_block\_shuffle},
\texttt{op\_block\_xor\_cbc}) — each with signature
\texttt{op\_xxx(text, encrypt, **kwargs)} and each claimed to be
self-contained and reversible.
\fussy

Six \emph{CipherLists} (CL1–CL6) each select exactly eight of these
primitives, in a fixed order, keyed to an input type inferred by a
heuristic auto-detector (\texttt{auto\_detect}): plain text, image,
video, document, numeric/data, and an ``experimental'' catch-all.
Encryption runs the eight steps forward; decryption runs the same eight
steps in reverse with \texttt{encrypt=False}. The identifier of the
CipherList used is returned alongside the ciphertext and is required for
decryption.

The Tkinter UI (\texttt{kod8\_ui.py}) wraps this engine, and — critically
for the threat model developed in Section~\ref{sec:threat} — persists
encrypted output to disk in a fixed text format:

\begin{lstlisting}[language=Python]
f.write(f"# KOD8  cipher_id: {cid}\n{enc}")
\end{lstlisting}

% =====================================================================
\section{Threat Model}
\label{sec:threat}
% =====================================================================

We evaluate KOD8 under three progressively weaker attacker assumptions,
each grounded in how the reference implementation actually behaves.

\begin{description}[leftmargin=1.4em, style=nextline]
  \item[\textbf{T1 — Source-available (white-box).}]
    \texttt{KOD8\_KEY} is a hardcoded module-level constant, not a
    user-supplied secret, a password-derived key, or a per-deployment
    value. Any party with access to \texttt{kod8\_engine.py} — including
    anyone who has ever received this repository, cloned it from a
    version-control host, or extracted it from a distributed build —
    possesses the key outright. Under Kerckhoffs's principle this would
    normally be an acceptable assumption \emph{provided the key itself
    remained secret}; here the key \emph{is} the source file. This is the
    single most severe practical weakness in the system and requires no
    cryptanalysis at all to exploit.
  \item[\textbf{T2 — Ciphertext-and-header-only, key rotated.}]
    A more charitable model assumes a deployer has changed
    \texttt{KOD8\_KEY} to a private eight-digit value not present in any
    public source tree, and that the attacker only observes
    \texttt{.kod8} files or wire traffic. Because the CipherList
    identifier is stored as a cleartext header (Section~\ref{sec:cid}),
    the attacker's search problem is immediately reduced to: recover one
    eight-digit key against one \emph{known} 8-step chain — not, as the
    ``six independent CipherLists'' design might suggest, a joint
    key-and-chain identification problem.
  \item[\textbf{T3 — Known-plaintext.}]
    The strongest assumption we exercise empirically in
    Section~\ref{sec:brute} is a single known plaintext/ciphertext pair,
    which is a routine assumption in classical cryptanalysis and
    realistic here: KOD8's own self-test harness, UI placeholder text,
    and predictable file headers (e.g.\ document boilerplate, JSON key
    names) are all plausible cribs.
\end{description}

The remainder of this report shows that KOD8 does not survive T2 or T3,
and trivially fails under T1.

% =====================================================================
\section{Per-Primitive Key-Dependence Analysis}
\label{sec:entropy}
% =====================================================================

The chain-safety documentation in \texttt{kod8\_engine.py} advertises
``6 CipherLists $\times$ 8 steps'' as the source of the scheme's
strength. We traced every primitive's use of \texttt{KOD8\_KEY} directly
in source and confirmed each classification by execution. Table~\ref{tab:entropy}
summarizes the result: of the primitives available to the engine, six
reference \texttt{KOD8\_KEY} in \emph{no way whatsoever}, and of the
remainder, several derive from only a small function of the key rather
than the full eight-digit value.

\begin{table}[h]
\centering
\small
\caption{Effective key-dependent state space per primitive, against a
nominal $10^{8}$-key eight-digit key. ``Zero'' means the operation is
fully determined by the algorithm alone once it is known — which,
per Kerckhoffs's principle, must be assumed.}
\label{tab:entropy}
\begin{tabular}{@{}lll@{}}
\toprule
\textbf{Primitive} & \textbf{Key dependence} & \textbf{Effective states} \\
\midrule
\texttt{op\_hex}          & none (encoding only)         & 1 (zero) \\
\texttt{op\_base64}       & none (encoding only)         & 1 (zero) \\
\texttt{op\_atbash\_tr}   & none                          & 1 (zero) \\
\texttt{op\_unicode\_shift} & none (delta hardcoded $\pm5$) & 1 (zero) \\
\texttt{op\_rail\_fence}  & none (rails hardcoded per CL) & 1 (zero) \\
\texttt{op\_block\_rotate} & none ($n$ hardcoded to 3)    & 1 (zero) \\
\texttt{op\_split\_reverse} & none                        & 1 (zero) \\
\texttt{op\_full\_reverse} & none                          & 1 (zero) \\
\texttt{op\_columnar}     & digit-sum mod 5 & \textbf{5} (columns $\in\{4..8\}$) \\
\texttt{op\_xor\_key}     & first digit only              & \textbf{10} \\
\texttt{op\_rolling\_xor} & first digit only (seed)       & \textbf{10} \\
\texttt{op\_sbox}         & weighted digit-sum ($\Sigma\, d_i \! \cdot\! (i{+}1)$) & \textbf{325} \\
\texttt{op\_block\_xor\_cbc} & IV = literal \texttt{KOD8\_KEY} string & fixed, non-random \\
\texttt{op\_block\_shuffle} & digits mod block count, deduplicated & $\le 8!$, usually far smaller \\
\texttt{op\_vigenere\_tr} & full key, period-8 repeating shift & $10^{8}$ (nominal) \\
\texttt{op\_base36}       & full key, period-8 repeating shift & $10^{8}$ (nominal) \\
\texttt{op\_keystream\_xor} & full key as LCG seed          & $10^{8}$ (nominal) \\
\bottomrule
\end{tabular}
\end{table}

Three findings from Table~\ref{tab:entropy} are worth stating explicitly.

\paragraph{(a) Six of seventeen primitives are not cryptographic at all.}
\texttt{op\_atbash\_tr}, \texttt{op\_unicode\_shift},
\texttt{op\_rail\_fence}, \texttt{op\_block\_rotate},
\texttt{op\_split\_reverse}, and \texttt{op\_full\_reverse} never inspect
\texttt{KOD8\_KEY}. We confirmed this by static inspection of every
function body:

\begin{lstlisting}[language=Python]
op_atbash_tr      -> references KOD8_KEY: False
op_unicode_shift  -> references KOD8_KEY: False
op_rail_fence     -> references KOD8_KEY: False
op_block_rotate   -> references KOD8_KEY: False
op_split_reverse  -> references KOD8_KEY: False
op_full_reverse   -> references KOD8_KEY: False
\end{lstlisting}

In CL1 (plain text), for example, four of the eight advertised layers —
\texttt{atbash\_tr}, \texttt{block\_rotate}, \texttt{unicode\_shift}, and
\texttt{full\_reverse} — fall into this category. Under Kerckhoffs's
assumption (the attacker knows the algorithm, only the key is secret),
these steps contribute \emph{no} confidentiality whatsoever: they are
fixed, publicly known, and instantly invertible by anyone who has read
this report or the source. Their function is purely to increase visual
and structural obfuscation, not adversarial work factor.

\paragraph{(b) The columnar-transposition parameter has only five possible values.}
\texttt{op\_columnar} derives its column count from
\texttt{(sum(int(d) for d in KOD8\_KEY) \% 5) + 4}. Since this is bounded
between 4 and 8 regardless of the underlying key, we confirmed empirically
that 100{,}000 randomly sampled eight-digit keys produce only the five
values $\{4,5,6,7,8\}$ — meaning an attacker can exhaustively enumerate
every possible columnar transposition used by CL4 in five trials.

\paragraph{(c) The S-box seed collapses an $10^{8}$ key into 325 states.}
The deterministic Fisher–Yates shuffle that produces the substitution
box used by CL2, CL5, and CL6 is seeded with a positionally weighted
digit sum,
\begin{lstlisting}[language=Python]
seed = sum(int(c) * (i + 1) for i, c in enumerate(KOD8_KEY))
\end{lstlisting} For an eight-digit key this expression
is bounded between 0 (all-zero key) and $9\times(1+2+\cdots+8)=324$,
i.e.\ \textbf{325 possible seed values, and therefore at most 325
distinct S-boxes}, independent of which of the $10^{8}$ nominal keys is
in use. We validated this by sampling 200{,}000 random eight-digit keys
and recording the resulting seed:

\begin{lstlisting}
Distinct S-box seeds observed in 200,000 random keys: 287
Theoretical maximum distinct seeds: 325
\end{lstlisting}

\noindent
This is close to the theoretical ceiling, as expected once the sample
size exceeds the state space by two orders of magnitude. Practically,
this means an attacker facing any CipherList that uses \texttt{op\_sbox}
never needs to search $10^{8}$ keys for that layer — a 325-entry
precomputed table of candidate S-boxes suffices, and each candidate can
be tested in microseconds.

% =====================================================================
\section{Fixed Initialization Vector in the CBC-Style Layer}
\label{sec:iv}
% =====================================================================

\texttt{op\_block\_xor\_cbc} is explicitly modeled on CBC-mode block
chaining, with each plaintext block XORed against the \emph{previous
ciphertext block}, seeded by an initialization vector:

\begin{lstlisting}[language=Python]
iv = (KOD8_KEY * ((block_size // len(KOD8_KEY)) + 1))[:block_size]
\end{lstlisting}

With the default \texttt{block\_size=8}, which is also the length of
\texttt{KOD8\_KEY}, this expression reduces to \texttt{iv == KOD8\_KEY}
verbatim — we confirmed this by direct execution:

\begin{lstlisting}
block_xor_cbc IV == KOD8_KEY literally: True '86247931'
\end{lstlisting}

This has two consequences that a real block cipher's CBC mode is
specifically designed to avoid:

\begin{enumerate}[leftmargin=1.6em]
\item \textbf{The IV is not random and is never varied per message.}
Every invocation of every CipherList that uses this primitive
(CL2, CL3) uses the identical IV. Combined with the determinism shown in
Section~\ref{sec:determinism}, this guarantees that identical inputs
always chain identically from the very first block.
\item \textbf{The IV is not distinct from the key.} In a textbook CBC
construction the IV need not be secret, but reusing key material
\emph{as} the IV needlessly collapses two independent security
parameters (confidentiality key, chaining diversifier) into one,
removing any benefit CBC-style chaining would otherwise provide against
identical-prefix leakage across messages encrypted under different
(hypothetical, rotated) keys.
\end{enumerate}

% =====================================================================
\section{Determinism and the Absence of Semantic Security}
\label{sec:determinism}
% =====================================================================

Because every layer of every CipherList is either unkeyed
(Section~\ref{sec:entropy}) or keyed only by the static
\texttt{KOD8\_KEY} with no per-message nonce, salt, or random IV
anywhere in the chain, KOD8 is fully deterministic: identical plaintext
always produces byte-identical ciphertext.

\begin{lstlisting}
Same plaintext -> same ciphertext (deterministic)? True
\end{lstlisting}

This was confirmed by encrypting the string \texttt{"THE MEETING IS AT
NOON"} twice, via two independent \texttt{Kod8} instances, and observing
identical output both times. Deterministic encryption fails the standard
notion of semantic security (IND-CPA): an observer who cannot decrypt a
single message can still detect message repetition, build a
frequency profile of recurring ciphertexts across a captured
communication stream, and mount straightforward replay or
dictionary-matching attacks — all without ever recovering the key.

% =====================================================================
\section{Absence of Integrity and Authentication}
% =====================================================================

No primitive in KOD8 produces or checks a message authentication code,
and \texttt{decrypt()} performs no integrity validation on its output.
Worse, several primitives are written to fail \emph{silently}. For
example:

\begin{lstlisting}[language=Python]
def op_hex(text: str, encrypt: bool, **_) -> str:
    if encrypt:
        return text.encode("utf-8").hex().upper()
    try:
        return bytes.fromhex(text).decode("utf-8")
    except Exception:
        return text
\end{lstlisting}

A malformed or tampered ciphertext does not raise an error here; it
silently falls through and returns the (still-encrypted, or
partially-mangled) input unchanged, which then continues through the
remaining seven reverse steps and is ultimately displayed to the user as
if it were valid plaintext. \texttt{op\_base64} exhibits the identical
pattern. This is a textbook malleability weakness: an attacker with
write access to the ciphertext channel can flip, truncate, or splice
bytes — particularly within the XOR-based layers, which are linear and
therefore preserve attacker-controlled deltas directly into the
recovered plaintext — with no cryptographic mechanism to detect the
tampering.

% =====================================================================
\section{Cleartext CipherList Disclosure}
\label{sec:cid}
% =====================================================================

KOD8's auto-detection mechanism (\texttt{auto\_detect}) is presented as
adding an extra layer of ambiguity: a ciphertext could, in principle,
have been produced by any of six different eight-step chains. In the
reference implementation this ambiguity does not survive contact with
the persistence layer. \texttt{kod8\_ui.py} writes every saved file as:

\begin{lstlisting}[language=Python]
f.write(f"# KOD8  cipher_id: {cid}\n{enc}")
\end{lstlisting}

and the decrypt path reads this header back out directly:

\begin{lstlisting}[language=Python]
if line.startswith("# KOD8  cipher_id:"):
    cid = line.split(":")[1].strip()
\end{lstlisting}

Every \texttt{.kod8} file therefore discloses, in cleartext, exactly
which of the six eight-step chains was used. This reduces the
attacker's problem from ``recover a key against an unknown one-of-six
chain'' to ``recover a key against one specific, fully known chain'' —
eliminating what would otherwise be the system's only source of
genuine combinatorial ambiguity beyond the key itself.

% =====================================================================
\section{Empirical Key Recovery}
\label{sec:brute}
% =====================================================================

To establish a concrete, conservative upper bound on attack cost, we
implemented a blind known-plaintext brute-force attack against CL1 (the
plain-text CipherList) that makes \emph{no} use of the entropy-collapse
findings in Section~\ref{sec:entropy} — i.e.\ it searches the full
nominal $10^{8}$ keyspace as if every digit of \texttt{KOD8\_KEY}
contributed independently to every step, which Table~\ref{tab:entropy}
shows is not actually the case.

\subsection{Method}

A plaintext (\texttt{"ATTACK AT DAWN TOMORROW"}) was encrypted once
under the reference key to obtain a known plaintext/ciphertext pair.
For each of $N$ randomly sampled eight-digit candidate keys, we
re-derived the dependent S-box exactly as the engine does at import
time, ran the full CL1 decryption chain, and compared the result to the
known plaintext. Per-trial wall-clock time was measured directly and
extrapolated linearly to the full $10^{8}$-key search.

\subsection{Results}

\begin{lstlisting}
2000 trials in 0.1676s -> 83.80 microseconds/trial
Found at trial 1000 (candidate = true key '86247931')
Estimated time for exhaustive 10^8 search
  (single core, unoptimized reference Python): 8379.7s ~= 2.33 hours
\end{lstlisting}

This is a \emph{worst-case, no-shortcuts} figure using an interpreted,
unoptimized reference implementation on a single CPU core. It is an
upper bound in every practical sense:

\begin{itemize}[leftmargin=1.6em]
\item \textbf{Parallelization.} The search is embarrassingly parallel
(each candidate key is independent). Across 32 cores, 2.33 hours falls
to under 4.4 minutes; on commodity cloud infrastructure or a GPU-hosted
implementation, recovery drops to seconds.
\item \textbf{Algorithmic shortcuts.} Section~\ref{sec:entropy} shows
that \texttt{op\_xor\_key} and \texttt{op\_rolling\_xor} depend only on
the first key digit (10 candidates), and that the S-box depends on only
325 candidate seeds rather than $10^{8}$ keys. A cryptanalyst exploiting
these facts — for example, solving \texttt{op\_xor\_key} directly via a
single-byte-XOR frequency or known-plaintext attack, and enumerating the
325-entry S-box table rather than the full key — reduces several
CipherLists' effective search space by five to six orders of magnitude,
to the point of near-instant recovery on a laptop.
\item \textbf{No implementation optimization was attempted.} The
benchmark uses pure-Python string operations with no vectorization,
caching, or compiled extensions; a systems-level reimplementation (C,
Rust, or even NumPy-vectorized Python) would be expected to improve
per-trial latency by one to two further orders of magnitude.
\end{itemize}

% =====================================================================
\section{Per-CipherList Summary}
% =====================================================================

Table~\ref{tab:cl-summary} aggregates the findings above per CipherList,
identifying the weakest key-dependent primitive in each chain — i.e.\
the layer a cryptanalyst would target first.

\begin{table}[h]
\centering
\small
\caption{Weakest link per CipherList and the resulting practical attack surface.}
\label{tab:cl-summary}
\begin{tabular}{@{}p{1.1cm}p{3.6cm}p{4.6cm}p{4.2cm}@{}}
\toprule
\textbf{ID} & \textbf{Input class} & \textbf{Weakest keyed primitive} & \textbf{Practical implication} \\
\midrule
CL1 & Plain text & \texttt{op\_xor\_key} (10 states) & Trivial single-byte-XOR solve; 4 of 8 steps unkeyed \\
CL2 & Image & \texttt{op\_sbox} (325 states) + fixed CBC IV & S-box table lookup; deterministic chaining from block 0 \\
CL3 & Video & fixed CBC IV (\texttt{op\_block\_xor\_cbc}) & Identical file prefixes leak identical ciphertext prefixes \\
CL4 & Document & \texttt{op\_columnar} (5 states) & Column count enumerable in 5 trials \\
CL5 & Numeric/data & \texttt{op\_sbox} (325 states) & S-box table lookup dominates over full-key search \\
CL6 & Experimental (fallback) & \texttt{op\_sbox} (325 states) & Same as CL2/CL5 \\
\bottomrule
\end{tabular}
\end{table}

% =====================================================================
\section{Discussion}
% =====================================================================

Taken together, these findings support a specific and fairly narrow
conclusion: \textbf{KOD8 is a well-constructed \emph{obfuscation}
layer and a legitimate example of classical product-cipher design, but
it does not meet the security expectations implied by its ``8-layer,
6-strategy'' framing, and it should not be relied upon to protect data
against a motivated adversary.}

This is not a criticism unique to KOD8; it is the expected outcome for
\emph{any} cipher built from XOR, substitution, and transposition
primitives over a small numeric key without a formal security proof,
a vetted key-derivation function, randomized nonces, and an
authenticated-encryption construction. Two observations frame this
constructively rather than dismissively:

\begin{enumerate}[leftmargin=1.6em]
\item \textbf{The stated design goal — resisting casual, automated,
pattern-based decoding (e.g.\ by an LLM skimming ciphertext for
recognizable structure) — is a materially different, and much weaker,
goal than confidentiality against a targeted cryptanalytic adversary.}
Under the former goal, several of the ``zero-entropy'' primitives in
Table~\ref{tab:entropy} (Atbash, transpositions, codepoint shifting)
are entirely reasonable: they destroy exactly the kind of surface-level
regularity that pattern-matching or naive statistical decoding would
key off, even though they add nothing against a human cryptanalyst who
reads the source. Evaluated against \emph{that} narrower goal, KOD8
likely performs its intended job.
\item \textbf{The gap between the two goals is exactly where the
current implementation is most exposed}, principally via T1 (hardcoded
key) and the cleartext CipherList header — both of which are
implementation choices, not inherent to the cipher-chain design, and
both are straightforward to remediate without redesigning the eight
primitive families.
\end{enumerate}

% =====================================================================
\section{Recommendations}
% =====================================================================

\begin{enumerate}[leftmargin=1.6em]
\item \textbf{Never treat KOD8 in its current form as confidentiality
protection for sensitive data.} If genuine confidentiality is required,
use a vetted AEAD construction (AES-256-GCM or ChaCha20-Poly1305) with a
key derived via a proper KDF (Argon2id, scrypt, or PBKDF2 with a modern
iteration count) from a user-supplied passphrase — not a hardcoded
constant.
\item \textbf{If KOD8 is retained for its stated obfuscation purpose},
document that scope explicitly in the code and UI so future users do
not mistake it for a security control, and stop persisting the
CipherList identifier in cleartext alongside the ciphertext — derive it
instead from a value only the key holder can compute (e.g.\ include it
inside the encrypted payload).
\item \textbf{Introduce a random, per-message nonce or IV} into any
retained CBC-style or keystream layer, transmitted alongside (not
derived from) the ciphertext, to eliminate the determinism identified
in Section~\ref{sec:determinism}.
\item \textbf{Add a MAC or AEAD tag} (even a simple HMAC-SHA256 over the
final ciphertext) so that tampering is detected rather than silently
decrypted into plausible-looking garbage, closing the failure mode in
Section~\ref{sec:iv}.
\item \textbf{Re-derive the S-box seed from the full key material}
(e.g.\ via a cryptographic hash of \texttt{KOD8\_KEY} rather than a
positional digit sum) to eliminate the 325-state collapse in
Section~\ref{sec:entropy}.
\end{enumerate}

% =====================================================================
\section{Conclusion}
% =====================================================================

We performed a source-grounded, empirically validated cryptanalysis of
KOD8. Static analysis of all seventeen primitives shows that only
three of the eight steps in a typical CipherList meaningfully depend on
the key, and that the two primitives most resembling ``real''
cryptographic components — the S-box and the CBC-style block XOR —
each contain implementation-level defects (a 325-state seed collapse and
a literal-key-as-IV reuse, respectively) that remove most of the
protection their design intends to provide. A conservative,
no-shortcuts, known-plaintext brute-force attack recovers the full
eight-digit key against the plain-text CipherList in under 2.4 hours on
a single CPU core with an unoptimized reference implementation, and the
structural weaknesses identified here reduce that further to seconds.
Combined with a hardcoded global key and a cleartext CipherList header
in the reference file format, KOD8 is best classified as a classical
obfuscation cipher suitable for resisting casual or automated
pattern-based inspection, not as a cryptographically secure scheme, and
it should be labeled and used accordingly.

\end{document}
